\section{Discussion and conclusion}

The thermal response separates operator cancellation from a ground-state
interference zero. Under the stated periodic symmetries, the doublet's
first-order probe can be canceled throughout the retained space; the
quartet's different bond-spin corrections prevent this. Its four-cell
ground state is exactly dark but supports an activated thermal response,
verified in the complete twenty-spin low band.
The exact second moment quantifies the obstruction for the prescribed
dipole--composite probe at general spin. Locality strengthens it to a
size-independent finite-frequency weight bound in the effective uniform
trace. Independent infinitesimal SU(2)-invariant controls inside each cell
cannot improve on uniform retuning: additional operator momenta add
nonnegative weight. Thus the residual is not an artifact of allowing only
one exchange knob.

This interpretation builds on established projection, tensor, and
effective-observable methods. In particular, higher-order scattering operators
are not introduced here for the first time~\cite{waldmann2005manyspin}, nor is
the correction of a cluster interference factor~\cite{mueller2003structure}.
The distinction between symmetry-fixed momentum dependence and
interference involving independent dynamical matrix elements is already
explicit in Ref.~\onlinecite{tabrizi2021selection}, Eqs.~(13)--(16).
Our result determines what remains after optimizing the physically
accessible coefficients of such amplitudes, rather than inferring a
residual simply from the presence of more than one operator.
Conditional excitation processes and thermal redistribution of spectral
weight are also established effects~\cite{exius2010vertex,honecker2016thermal}.
Even the thermal lifting of an accidental scattering zero has an
experimental precedent in magnetovibrational interference~\cite{raymond2018interference}.
The distinction examined here is whether a physical exchange retuning can
remove the corrected response simultaneously across the thermally occupied
low-energy manifold. The surviving vector acts on neighboring effective
spins; it is not a direct rank-two probe or a new excitation species.
The local representation constrains this operator, while ground-state
accessibility determines whether its independence is already visible at
zero temperature.

Several qualifications determine the scope of that distinction. The local
spin permits a tensor but does not ensure that it acts independently on the
many-body ground state, as the four-cell periodic exception demonstrates.
Open endpoints create a separate one-spin residual. The microscopic
same-geometry test is an odd frustrated ring; the contrast survives an even
ring at the same roots, but that check uses the leading effective model.
Converged twelve-cell periodic data rule out a purely endpoint explanation
at the tested quartet root. The finite-temperature commutator proof removes
a dependence on particular even chain lengths, but not the fixed-size
weak-coupling assumption. The weight bound in
Eq.~\eqref{eq:quartet-window-bound} applies to the effective uniform trace,
with that limit taken before the size limit. It is not a thermodynamic
lower bound at arbitrary fixed scaled temperature and does not establish
uniform control of the microscopic expansion with size. The general-spin
results also presuppose a specified nonzero composite probe coefficient.
The local-control extension does not cover anisotropy, changes of local
irrep, or modifications of the leading intercell bond ratios.
In the four-cell microscopic check, the fourth-order
ground-state background eventually dominates the activated response at
fixed nonzero coupling. Higher-order level splittings also generate
quasielastic weight absent from the leading finite-frequency coefficient.
The order of the weak-coupling, low-temperature, and frequency-window
limits therefore matters. Furthermore,
finite momentum acceptance can overwhelm the ideal doublet signal with
ordinary form-factor weight. Window-optimized exchange removes the mean
form-factor shift but leaves its variance; the required shrinking-window
scales depend on the order of the selected momentum zero.
There is no material-specific demonstration of detectability.
The doublet's observed fourth-order behavior along the first-order
trajectory is not an unavoidable floor. Further retuning numerically
cancels the central response in the two-cell example, whose retained
space contains only one singlet and one triplet. A complete second-order
probe and second-order retuning for general chains have not been analyzed.

Within these limits, the pentamer chain gives a quantitative distinction
between an operator-level cancellation and a state-dependent dark point.
Its value is in identifying which inference can be drawn from weak
scattering: a dark zero-temperature spectrum does not by itself imply a
thermal selection rule or absence of the corresponding excitations.
Using the distinction experimentally requires control of momentum
acceptance and nonleading backgrounds in addition to the effective
Hamiltonian and local spin representation.
